\newcommand{\ModuleID}{EL-I.7}
\newcommand{\ModuleTitle}{\Latin{Sum} and \Latin{Possum}}
\newcommand{\ModuleStage}{Stage I — Foundations}
\newcommand{\ModuleUnit}{I.7}
\newcommand{\ModuleOutcome}{Recognize the six indicative tenses of \Latin{sum}, use \Latin{possum} with its infinitive, and distinguish predicate complements from objects}
\newcommand{\ModulePrerequisites}{EL-I.6 mastery: trace noun--adjective agreement and classify adjective use}
\newcommand{\ModuleSessionOne}{Triage copula, predicate, and complement structures across the full indicative of \Latin{sum}, then complete Worksheet I.26 as the guided all-tense application.}
\newcommand{\ModuleSessionTwo}{Retrieve \Latin{possum} and its complementary-infinitive pattern, complete Worksheet I.27 independently, and pass the enrichment exit ticket without assigning accusative case to a predicate complement.}
\newcommand{\ModulePrintedPractice}{Worksheets I.26 (Variant B) and I.27 (Variant C), with terminal keys}
\newcommand{\ModuleExtensionPractice}{Worksheets I.25 and I.28 remain optional remediation or delayed-recall variants}
\newcommand{\ModuleNext}{EL-I.8, The Present Active System}
\newcommand{\ModuleMastery}{Write at least 34 of 36 displayed \Latin{sum} forms correctly, parse at least 18 of 20 \Latin{sum}/\Latin{possum} forms, and analyze five linking or ability clauses without confusing predicate, object, or infinitive complement.}
\newcommand{\ModuleDelayNote}{}
\newcommand{\ModuleReferenceFocus}{%
  \begin{itemize}
    \item the present, imperfect, future, perfect, pluperfect, and future-perfect indicative of \Latin{sum};
    \item nominative predicate nouns and adjectives after a copula;
    \item ordinary indicative forms and perfect stems of \Latin{possum}; and
    \item complementary infinitives and visibly bracketed supplied copulas.
  \end{itemize}}
\newcommand{\ModuleContrast}{\Latin{Sum} links a subject to a nominative predicate complement; it does not transfer action to an accusative object. \Latin{Possum} normally requires an infinitive that completes its meaning, not a predicate nominative.}
\newcommand{\ModuleLessonSource}{\CourseRoot/shared/content/volume1/all}
\newcommand{\ModuleExerciseSource}{\CourseRoot/shared/exercises/volume1/all}
\newcommand{\ModuleWorkedBank}{I}
\newcommand{\ModuleExerciseBank}{I}
\newcommand{\ModuleWorksheetVariants}{B,C}
\newcommand{\ModuleScopeNote}{This packet teaches the indicative of \Latin{sum} and operational \Latin{possum}; complete subjunctive morphology remains in Stage II.}
\newcommand{\ModuleEnrichment}{%
  \SubmoduleExtension
    {The indicative of \Latin{sum}}
    {A copula/predicate/complement triage begins with the verb's function rather than an English word order.}
    {For the project-composed clauses \Latin{Maria magistra est}, \Latin{Discipuli parati erant}, and \Latin{Pax erit donum}, box the copula, underline the subject, and double-underline every predicate complement. Then change only the tense.}
    {State the case of both subject and predicate complement and explain why neither is a direct object.}
    {In each project-composed clause, the subject and predicate noun or adjective are nominative; \Latin{est}, \Latin{erant}, and \Latin{erit} link rather than transfer action. A tense change alters the copula while preserving the nominal relation unless another coordinate is deliberately changed.}
  \SubmoduleExtension
    {The compound \Latin{possum}}
    {The ability verb and its infinitive form one verbal center while retaining distinct finite and non-finite jobs.}
    {For the project-composed clause \Latin{Discipulus potest legere}, label person, number, tense, and mood on \Latin{potest}; label \Latin{legere} as complementary infinitive; then change the subject to plural and the time to past ability.}
    {Explain which verb anchors the clause and which form supplies the activity made possible.}
    {The transformed clause may be \Latin{Discipuli poterant legere}. Finite \Latin{potest/poterant} anchors person, number, tense, and mood; \Latin{legere} is non-finite and completes the idea of ability.}
  \SubmoduleExtension
    {Handwritten study and controlled composition}
    {An independent exit ticket tests whether paradigms support clause analysis instead of remaining isolated recitation.}
    {Without notes: parse \Latin{fueramus, erunt, potuistis}; supply nominative predicate forms for a plural subject; build one project-composed ability clause; and bracket an omitted copula in a verbless project-composed acclamation.}
    {Write one sentence distinguishing object, predicate complement, and complementary infinitive.}
    {Expected parses are first-person plural pluperfect of \Latin{sum}, third-person plural future of \Latin{sum}, and second-person plural perfect of \Latin{possum}. Predicate forms must agree with the plural subject; the ability clause needs finite \Latin{possum} plus an infinitive; a supplied copula must be visibly bracketed. The final distinction must identify accusative target, nominative completion, and non-finite verbal completion respectively.}
}
